Source : DVF 2023-2024 (prix médian au m² de la commune) ; CGI art. 1594 D et 1584 (droits) ; modèle du bloc (taxe foncière) ; note8p/calculs/foncier_calculs.py.\par Droits de mutation au taux permanent (4,50 % département, 1,20 % commune, frais d'assiette de l'État), ceux que paient en 2026 les primo-accédants et tous les acheteurs à partir du 1er avril 2028 ; les autres acheteurs paient 6,32 % jusqu'au 31 mars 2028 dans ces trois départements. Hausse de taxe foncière : enveloppe de 15 Md€, en régime permanent (atteinte en trois paliers), proportionnelle à la valeur dans la commune.
